\chapter{Evaluation Harness and Static Outputs}
\label{app:evaluation-harness}

This appendix provides total forensic transparency regarding the prototype's evaluation metrics. It explicitly dissects the origin of the reported Precision, Recall, and F1 scores, proving their status as static architectural stubs, and establishes the mathematical blueprints required to compute empirical metrics in the future.

\section{Static Output Generation Forensics}
\label{app-sec:eval-forensics}
The Vzeya cinematic frontend prominently displays the following metrics:
\begin{itemize}
    \item \textbf{Precision:} $0.9280$
    \item \textbf{Recall:} $0.4598$
    \item \textbf{F1 Score:} $0.6149$
\end{itemize}

\textbf{Forensic Origin:} These values are explicitly \textbf{NOT} the result of dynamic Jaro-Winkler computation over the \texttt{postgres-truth} vault. They are generated by the \texttt{generator/src/evaluator/\_\_main\_\_.py} script. Inspection of this source file reveals that these numbers are formatted string literals printed to standard output. 

\textbf{Scientific Implication:} The evaluator module functions purely as an architectural CI/CD test (Class D). It proves that an external process can connect to the database, simulate a reporting cycle, and output a JSON format readable by the frontend. Using these values to claim empirical accuracy for Wolverine's entity resolution is mathematically invalid.

\section{Dynamic Evaluation Blueprint (Future Work)}
\label{app-sec:eval-blueprint}
To generate valid (Class A) metrics in a future iteration (\Cref{ch:future-work}), the evaluator must be rewritten to compute a dynamic Confusion Matrix against the 39,605 generated ground-truth edges.

\subsection{Confusion Matrix Definitions}
\begin{itemize}
    \item \textbf{True Positive (TP):} A predicted edge ($S \ge 0.70$) that exactly matches a ground-truth edge.
    \item \textbf{False Positive (FP):} A predicted edge ($S \ge 0.70$) that does \textit{not} exist in the ground truth (False Attribution).
    \item \textbf{False Negative (FN):} A ground-truth edge that the system failed to predict ($S < 0.70$).
\end{itemize}

\subsection{Metric Equations}
\begin{equation}
    \text{Precision (P)} = \frac{TP}{TP + FP}
\end{equation}
\begin{equation}
    \text{Recall (R)} = \frac{TP}{TP + FN}
\end{equation}
\begin{equation}
    \text{F1 Score} = 2 \times \frac{P \times R}{P + R}
\end{equation}

\textit{Note: The equations above are methodological blueprints. The prototype's static metrics must not be inserted into these formulas.}

\section{Threshold Evaluation Blueprint}
\label{app-sec:eval-threshold}
The system currently utilizes static decision thresholds:
\begin{itemize}
    \item $S < 0.70$: Discarded / Unlinked.
    \item $0.70 \le S < 0.92$: Flagged for Human Analyst Review.
    \item $S \ge 0.92$: Automatically Linked (High Confidence).
\end{itemize}
Future evaluation must implement a threshold sweep (e.g., $S \in [0.50, 0.99]$) to plot empirical Precision-Recall (PR) and Receiver Operating Characteristic (ROC) curves. This calibration is necessary to quantify the trade-off between missing threat actor linkages (FN) and generating overwhelming analyst alert fatigue (FP).

\section{Scalability Benchmark Blueprint}
\label{app-sec:eval-scalability}
The BFS projection is theoretically bounded at $O(|V| + |E|)$. However, this asymptotic complexity is shielded by the hard limits $d \le 4, N \le 500$ implemented in the V8 engine. A future scalability benchmark (Class E) must disable these hard limits and measure query latency (ms) and peak memory consumption (MB) against a disk-backed graph database as $N$ scales from $10^3$ to $10^6$ nodes.

\section{Reproducibility Procedure}
\label{app-sec:eval-repro}
To reproduce the static evaluator stub behavior:
\begin{enumerate}
    \item Ensure the Docker environment is running (\Cref{app:synthetic-generator}).
    \item \textbf{Command:} \texttt{docker-compose run generator python -m evaluator}
    \item \textbf{Expected Output:} The console will print the static JSON payload containing the $0.9280$, $0.4598$, and $0.6149$ values.
\end{enumerate}
